\section{Threats to Validity}
\label{sec:threats}

\paragraph{Internal threats.} A harness defect can look like an engine defect.
During development our own adapter once read an error message as a result row
(Section~\ref{sec:impl}). We guarded against this as follows:
\begin{itemize}
  \item every reported statement was re-run by hand on a fresh process, on the
        current upstream main branch, before it was filed;
  \item each defect was confirmed on two builds of the same source tree that
        differ only by our patches, so the patch explains the change;
  \item a test engine with planted defects checks that the oracle fires on a real
        defect and stays silent on a correct engine and on invalid SQL;
  \item every seed in Table~\ref{tab:defects} was run again in a single case and
        reproduced the signature claimed.
\end{itemize}
What remains: each run in Table~\ref{tab:ablation} was done once. Seeds make a
run repeatable on the same code, so the counts there are exact for that code.
The first seed at which a source reaches a defect is one sample and not a rate.
Timings come from one laptop and will differ elsewhere. The claims rest on exit
status and error text, which do not vary between our runs.

\paragraph{External threats.} We tested one engine, mined one other engine, and
all four defects are in one family of string functions in one file. Three share a
cause. Nothing here shows that the method finds defects in another engine, and
Section~\ref{sec:discussion} gives two reasons to doubt that it transfers cheaply:
it needs a rich, labelled tracker, and most mined reports do not run on a
different dialect. The comparison in Table~\ref{tab:ablation} measures \tool as
it stands now. The generator gained \sql{lpad}, \sql{rpad} and \sql{repeat} only
after the first defects were filed, so the table does not say which source would
have found them first at that time.

\paragraph{Construct threats.} We measure the blind spot through one table of
functions, one literal policy and a text search of one source tree. That is a
sound lower bound: a generator cannot emit a function it does not name or a value
outside its configured range. It is not a proof that no other code path in the
project reaches the same statement. We also count defects by signature. A
signature is a hash of the error, so two causes can share one, as defects 1
and~2 do. For that reason we decided the count by reading the code, not from
signatures. The blind-spot argument itself was formed after the first defect was
found, as an explanation of why it had survived. We present it as a hypothesis
that the evidence supports, not as the reason the first defect was found.
